% assets/w03-6dof.svg — 6자유도: 힘·속도·위치가 어떻게 짝지어지는가.
%
%   bash _tools/tikz2svg.sh figures/src/w01-6dof.tex assets/w03-6dof.svg
%
% 3차원은 눈대중으로 기울이지 않는다 . TikZ 에 축 셋의
% 화면 사영을 직접 주면 좌표 (x,y,z) 를 그대로 쓸 수 있고, 사영은 TikZ 가 한다.
% 사영은 대학원 과목 자료의 같은 그림 가 쓰던 것과 같다:
%
%   P(x) = ( 0.87687, -0.21477)   선수  — 오른쪽 위 (화면 아래 성분이 음수)
%   P(y) = ( 0.48072,  0.39176)   우현  — 오른쪽 아래
%   P(z) = ( 0.00000,  0.89465)   아래
%
% TikZ 의 y 는 위가 양수이므로 두 번째 성분의 부호를 뒤집어 넣는다.
%
% 모멘트 호는 그 축에 수직인 평면 위의 원이다. 3차원 좌표로 원을 매개하면
% 사영된 타원과 그 중심이 축선 위에 오는 것이 저절로 보장된다 — 이것이
% 손으로 타원을 놓지 않는 이유다.
\documentclass[border=5pt]{standalone}
% capstone-style.tex — 이 볼트의 TikZ 그림이 공유하는 스타일
%
% 쓰는 법:  % capstone-style.tex — 이 볼트의 TikZ 그림이 공유하는 스타일
%
% 쓰는 법:  % capstone-style.tex — 이 볼트의 TikZ 그림이 공유하는 스타일
%
% 쓰는 법:  \input{capstone-style}   (assets/src/ 안에서)
% 빌드:     bash _tools/tikz2svg.sh assets/src/w02-node-topic.tex assets/w02-node-topic.svg
%
% 왜 TikZ 인가
%   손으로 SVG 좌표를 쓰면 화살촉·선 굵기·글자 크기가 그림마다 미묘하게 달라진다.
%   그것이 "자료가 어설퍼 보인다"의 정체다. 스타일을 한 번만 정하고 상대좌표로
%   배치하면 좌표를 고쳐도 그림이 무너지지 않는다.
%   수식은 본문과 같은 TeX 글꼴로 나온다.
%
% 한글
%   ko.TeX(DVI 경로)를 쓴다. XeLaTeX 은 TikZ 그래픽을 PDF special 로 내보내는데
%   dvisvgm 이 그것을 받으려면 Ghostscript < 10.01 이 필요하다. 이 PC 에는
%   10.07 이 깔려 있어 그 경로가 막힌다. latex -> DVI -> dvisvgm 이 유일하게
%   동작하는 조합이다.

\usepackage[hangul]{kotex}
\usepackage{tikz}
\usepackage{amsmath}
\usepackage{xcolor}
\usetikzlibrary{arrows.meta, positioning, calc, fit, backgrounds, decorations.pathmorphing}

% ---- 색 --------------------------------------------------------------------
% 강조는 하나만. 나머지는 검정. 흑백 인쇄에서도 읽혀야 한다.
\definecolor{ink}{HTML}{1A1A1A}
\definecolor{soft}{HTML}{666666}
\definecolor{accent}{HTML}{7C3AED}
\definecolor{warn}{HTML}{B91C1C}

% 계층 색 — Simulink 모델의 블록 색과 같은 규칙을 쓴다
\definecolor{guid}{HTML}{1D4ED8}    % 유도
\definecolor{ctrl}{HTML}{EA580C}    % 제어
\definecolor{thr}{HTML}{D97706}     % 추진기
\definecolor{plant}{HTML}{16A34A}   % 운동모델
\definecolor{nav}{HTML}{7C3AED}     % 항법 · 신호처리
\definecolor{pathc}{HTML}{0D9488}   % 계획 경로 · 항적
\definecolor{errc}{HTML}{C2410C}    % 오차
\definecolor{flow}{HTML}{0369A1}    % 조류 · 바람

% ---- 선과 화살촉 -----------------------------------------------------------
\tikzset{
  sig/.style   = {ink, line width=0.5pt, -{Stealth[length=4pt, width=3pt]}},
  wire/.style  = {ink, line width=0.5pt},
  fb/.style    = {ink, line width=0.5pt, densely dashed,
                  -{Stealth[length=4pt, width=3pt]}},
  asig/.style  = {accent, line width=0.5pt, densely dashed,
                  -{Stealth[length=4pt, width=3pt]}},
  % 블록 — 흰 바탕, 직각, 검은 테두리. 색은 테두리에만 준다
  blk/.style   = {draw=ink, line width=0.55pt, fill=white,
                  minimum width=18mm, minimum height=8.5mm, inner sep=2pt, font=\small},
  gblk/.style  = {blk, draw=guid},
  cblk/.style  = {blk, draw=ctrl},
  tblk/.style  = {blk, draw=thr},
  pblk/.style  = {blk, draw=plant},
  nblk/.style  = {blk, draw=nav},
  % 합산점 · 분기점
  sum/.style   = {draw=ink, line width=0.55pt, fill=white, circle,
                  inner sep=0pt, minimum size=4.6mm, font=\scriptsize},
  dot/.style   = {ink, fill=ink, circle, inner sep=0pt, minimum size=1.5mm},
  % 글자
  sname/.style = {font=\small, inner sep=1.5pt},
  note/.style  = {font=\scriptsize, soft, inner sep=1.5pt},
  wnote/.style = {font=\scriptsize, warn, inner sep=1.5pt},
  anote/.style = {font=\scriptsize, accent, inner sep=1.5pt},
  axis/.style  = {ink, line width=0.5pt},
  tick/.style  = {font=\scriptsize, soft},
  % 묶음 상자 — 서브시스템 표시
  grp/.style   = {draw=soft, line width=0.4pt, densely dotted, rounded corners=1mm,
                  inner sep=3mm},
}
   (assets/src/ 안에서)
% 빌드:     bash _tools/tikz2svg.sh assets/src/w02-node-topic.tex assets/w02-node-topic.svg
%
% 왜 TikZ 인가
%   손으로 SVG 좌표를 쓰면 화살촉·선 굵기·글자 크기가 그림마다 미묘하게 달라진다.
%   그것이 "자료가 어설퍼 보인다"의 정체다. 스타일을 한 번만 정하고 상대좌표로
%   배치하면 좌표를 고쳐도 그림이 무너지지 않는다.
%   수식은 본문과 같은 TeX 글꼴로 나온다.
%
% 한글
%   ko.TeX(DVI 경로)를 쓴다. XeLaTeX 은 TikZ 그래픽을 PDF special 로 내보내는데
%   dvisvgm 이 그것을 받으려면 Ghostscript < 10.01 이 필요하다. 이 PC 에는
%   10.07 이 깔려 있어 그 경로가 막힌다. latex -> DVI -> dvisvgm 이 유일하게
%   동작하는 조합이다.

\usepackage[hangul]{kotex}
\usepackage{tikz}
\usepackage{amsmath}
\usepackage{xcolor}
\usetikzlibrary{arrows.meta, positioning, calc, fit, backgrounds, decorations.pathmorphing}

% ---- 색 --------------------------------------------------------------------
% 강조는 하나만. 나머지는 검정. 흑백 인쇄에서도 읽혀야 한다.
\definecolor{ink}{HTML}{1A1A1A}
\definecolor{soft}{HTML}{666666}
\definecolor{accent}{HTML}{7C3AED}
\definecolor{warn}{HTML}{B91C1C}

% 계층 색 — Simulink 모델의 블록 색과 같은 규칙을 쓴다
\definecolor{guid}{HTML}{1D4ED8}    % 유도
\definecolor{ctrl}{HTML}{EA580C}    % 제어
\definecolor{thr}{HTML}{D97706}     % 추진기
\definecolor{plant}{HTML}{16A34A}   % 운동모델
\definecolor{nav}{HTML}{7C3AED}     % 항법 · 신호처리
\definecolor{pathc}{HTML}{0D9488}   % 계획 경로 · 항적
\definecolor{errc}{HTML}{C2410C}    % 오차
\definecolor{flow}{HTML}{0369A1}    % 조류 · 바람

% ---- 선과 화살촉 -----------------------------------------------------------
\tikzset{
  sig/.style   = {ink, line width=0.5pt, -{Stealth[length=4pt, width=3pt]}},
  wire/.style  = {ink, line width=0.5pt},
  fb/.style    = {ink, line width=0.5pt, densely dashed,
                  -{Stealth[length=4pt, width=3pt]}},
  asig/.style  = {accent, line width=0.5pt, densely dashed,
                  -{Stealth[length=4pt, width=3pt]}},
  % 블록 — 흰 바탕, 직각, 검은 테두리. 색은 테두리에만 준다
  blk/.style   = {draw=ink, line width=0.55pt, fill=white,
                  minimum width=18mm, minimum height=8.5mm, inner sep=2pt, font=\small},
  gblk/.style  = {blk, draw=guid},
  cblk/.style  = {blk, draw=ctrl},
  tblk/.style  = {blk, draw=thr},
  pblk/.style  = {blk, draw=plant},
  nblk/.style  = {blk, draw=nav},
  % 합산점 · 분기점
  sum/.style   = {draw=ink, line width=0.55pt, fill=white, circle,
                  inner sep=0pt, minimum size=4.6mm, font=\scriptsize},
  dot/.style   = {ink, fill=ink, circle, inner sep=0pt, minimum size=1.5mm},
  % 글자
  sname/.style = {font=\small, inner sep=1.5pt},
  note/.style  = {font=\scriptsize, soft, inner sep=1.5pt},
  wnote/.style = {font=\scriptsize, warn, inner sep=1.5pt},
  anote/.style = {font=\scriptsize, accent, inner sep=1.5pt},
  axis/.style  = {ink, line width=0.5pt},
  tick/.style  = {font=\scriptsize, soft},
  % 묶음 상자 — 서브시스템 표시
  grp/.style   = {draw=soft, line width=0.4pt, densely dotted, rounded corners=1mm,
                  inner sep=3mm},
}
   (assets/src/ 안에서)
% 빌드:     bash _tools/tikz2svg.sh assets/src/w02-node-topic.tex assets/w02-node-topic.svg
%
% 왜 TikZ 인가
%   손으로 SVG 좌표를 쓰면 화살촉·선 굵기·글자 크기가 그림마다 미묘하게 달라진다.
%   그것이 "자료가 어설퍼 보인다"의 정체다. 스타일을 한 번만 정하고 상대좌표로
%   배치하면 좌표를 고쳐도 그림이 무너지지 않는다.
%   수식은 본문과 같은 TeX 글꼴로 나온다.
%
% 한글
%   ko.TeX(DVI 경로)를 쓴다. XeLaTeX 은 TikZ 그래픽을 PDF special 로 내보내는데
%   dvisvgm 이 그것을 받으려면 Ghostscript < 10.01 이 필요하다. 이 PC 에는
%   10.07 이 깔려 있어 그 경로가 막힌다. latex -> DVI -> dvisvgm 이 유일하게
%   동작하는 조합이다.

\usepackage[hangul]{kotex}
\usepackage{tikz}
\usepackage{amsmath}
\usepackage{xcolor}
\usetikzlibrary{arrows.meta, positioning, calc, fit, backgrounds, decorations.pathmorphing}

% ---- 색 --------------------------------------------------------------------
% 강조는 하나만. 나머지는 검정. 흑백 인쇄에서도 읽혀야 한다.
\definecolor{ink}{HTML}{1A1A1A}
\definecolor{soft}{HTML}{666666}
\definecolor{accent}{HTML}{7C3AED}
\definecolor{warn}{HTML}{B91C1C}

% 계층 색 — Simulink 모델의 블록 색과 같은 규칙을 쓴다
\definecolor{guid}{HTML}{1D4ED8}    % 유도
\definecolor{ctrl}{HTML}{EA580C}    % 제어
\definecolor{thr}{HTML}{D97706}     % 추진기
\definecolor{plant}{HTML}{16A34A}   % 운동모델
\definecolor{nav}{HTML}{7C3AED}     % 항법 · 신호처리
\definecolor{pathc}{HTML}{0D9488}   % 계획 경로 · 항적
\definecolor{errc}{HTML}{C2410C}    % 오차
\definecolor{flow}{HTML}{0369A1}    % 조류 · 바람

% ---- 선과 화살촉 -----------------------------------------------------------
\tikzset{
  sig/.style   = {ink, line width=0.5pt, -{Stealth[length=4pt, width=3pt]}},
  wire/.style  = {ink, line width=0.5pt},
  fb/.style    = {ink, line width=0.5pt, densely dashed,
                  -{Stealth[length=4pt, width=3pt]}},
  asig/.style  = {accent, line width=0.5pt, densely dashed,
                  -{Stealth[length=4pt, width=3pt]}},
  % 블록 — 흰 바탕, 직각, 검은 테두리. 색은 테두리에만 준다
  blk/.style   = {draw=ink, line width=0.55pt, fill=white,
                  minimum width=18mm, minimum height=8.5mm, inner sep=2pt, font=\small},
  gblk/.style  = {blk, draw=guid},
  cblk/.style  = {blk, draw=ctrl},
  tblk/.style  = {blk, draw=thr},
  pblk/.style  = {blk, draw=plant},
  nblk/.style  = {blk, draw=nav},
  % 합산점 · 분기점
  sum/.style   = {draw=ink, line width=0.55pt, fill=white, circle,
                  inner sep=0pt, minimum size=4.6mm, font=\scriptsize},
  dot/.style   = {ink, fill=ink, circle, inner sep=0pt, minimum size=1.5mm},
  % 글자
  sname/.style = {font=\small, inner sep=1.5pt},
  note/.style  = {font=\scriptsize, soft, inner sep=1.5pt},
  wnote/.style = {font=\scriptsize, warn, inner sep=1.5pt},
  anote/.style = {font=\scriptsize, accent, inner sep=1.5pt},
  axis/.style  = {ink, line width=0.5pt},
  tick/.style  = {font=\scriptsize, soft},
  % 묶음 상자 — 서브시스템 표시
  grp/.style   = {draw=soft, line width=0.4pt, densely dotted, rounded corners=1mm,
                  inner sep=3mm},
}


\begin{document}
\begin{tikzpicture}[
    x={(0.87687cm,0.21477cm)},
    y={(0.48072cm,-0.39176cm)},
    z={(0cm,-0.89465cm)},
    scale=3.4]

% ---- 축 --------------------------------------------------------------------
\draw[axis, -{Stealth[length=4.5pt,width=3.4pt]}] (0,0,0) -- (1.5,0,0);
\draw[axis, -{Stealth[length=4.5pt,width=3.4pt]}] (0,0,0) -- (0,1.5,0);
\draw[axis, -{Stealth[length=4.5pt,width=3.4pt]}] (0,0,0) -- (0,0,1.5);
\node[font=\scriptsize, soft] at (1.72,0,0)  {$x_b$};
\node[font=\scriptsize, soft] at (0,1.66,0)  {$y_b$};
\node[font=\scriptsize, soft] at (0,0,1.66)  {$z_b$};
\fill[ink] (0,0,0) circle (0.6pt);
\node[font=\scriptsize, soft, anchor=east] at (-0.12,0,0) {$o_b$};

% ---- 갑판 (선수가 +x) -------------------------------------------------------
\fill[ink, opacity=0.10]
      (0.62,0,0) -- (0.34,0.26,0) -- (-0.53,0.26,0) -- (-0.62,0.15,0)
   -- (-0.62,-0.15,0) -- (-0.53,-0.26,0) -- (0.34,-0.26,0) -- cycle;
\draw[soft, line width=0.3pt]
      (0.62,0,0) -- (0.34,0.26,0) -- (-0.53,0.26,0) -- (-0.62,0.15,0)
   -- (-0.62,-0.15,0) -- (-0.53,-0.26,0) -- (0.34,-0.26,0) -- cycle;

% ---- 힘 셋 (직선 화살표) ----------------------------------------------------
\draw[thr, line width=1.0pt, -{Stealth[length=5pt,width=3.6pt]}]
      (0.75,0,0) -- (1.28,0,0);
\node[thr, font=\scriptsize] at (1.14,0.22,0) {$X$ — 전후};
\draw[thr, line width=1.0pt, -{Stealth[length=5pt,width=3.6pt]}]
      (0,0.45,0) -- (0,1.28,0);
\node[thr, font=\scriptsize, anchor=north] at (0,1.20,0.22) {$Y$ — 좌우};
\draw[thr, line width=1.0pt, -{Stealth[length=5pt,width=3.6pt]}]
      (0,0,0.30) -- (0,0,1.28);
\node[thr, font=\scriptsize, anchor=west] at (0.10,0,1.24) {$Z$ — 상하};

% ---- 모멘트 셋 (축에 수직인 평면 위의 원) ----------------------------------
% x 축 둘레 : y -> z 방향으로 도는 원, 중심은 x 축 위
\draw[accent, line width=0.9pt, -{Stealth[length=4.5pt,width=3.4pt]}]
      plot[domain=40:340, samples=60, variable=\t]
      ({2.05}, {0.30*cos(\t)}, {0.30*sin(\t)});
\node[accent, font=\scriptsize, anchor=south] at (2.05,0,-0.42) {$K$ — 횡경사};

% y 축 둘레 : z -> x
\draw[accent, line width=0.9pt, -{Stealth[length=4.5pt,width=3.4pt]}]
      plot[domain=40:340, samples=60, variable=\t]
      ({0.30*sin(\t)}, {1.85}, {0.30*cos(\t)});
\node[accent, font=\scriptsize, anchor=south] at (0,1.85,-0.42) {$M$ — 종경사};

% z 축 둘레 : x -> y
\draw[accent, line width=0.9pt, -{Stealth[length=4.5pt,width=3.4pt]}]
      plot[domain=40:340, samples=60, variable=\t]
      ({0.30*cos(\t)}, {0.30*sin(\t)}, {1.90});
\node[accent, font=\scriptsize, anchor=west] at (0,0.52,1.90) {$N$ — 선수각};

% 축선을 호 중심까지 점선으로 — "이 축 둘레" 임을 보이게 한다
\draw[soft, line width=0.3pt, densely dotted] (1.5,0,0) -- (2.05,0,0);
\draw[soft, line width=0.3pt, densely dotted] (0,1.5,0) -- (0,1.85,0);
\draw[soft, line width=0.3pt, densely dotted] (0,0,1.5) -- (0,0,1.90);
\fill[accent] (2.05,0,0) circle (0.4pt);
\fill[accent] (0,1.85,0) circle (0.4pt);
\fill[accent] (0,0,1.90) circle (0.4pt);

\end{tikzpicture}

\hspace{4mm}
% ===================== 오른쪽 : 표 =========================================
\begin{tikzpicture}[x=1mm, y=1mm]
\node[ink, font=\small\bfseries, anchor=north west] at (0,64)
     {선박의 6자유도};
\node[anchor=north west, align=left, text width=78mm, font=\scriptsize, soft] at (0,58) {%
  한 줄이 곧 하나의 인과다. \textbf{가로로} 읽는다 —
  힘을 주면 속도가 생기고, 그 속도를 적분하면 위치가 된다.};

\node[anchor=north west, font=\scriptsize] at (0,44) {%
  \begin{tabular}{@{}lllll@{}}
    \textbf{\#} & \textbf{자유도} & \textbf{힘 · 모멘트} & \textbf{속도} & \textbf{위치 · 자세}\\[2pt]
    1 & 전후 (surge) & $X$ & $u$ & $x$\\
    2 & 좌우 (sway)  & $Y$ & $v$ & $y$\\
    3 & 상하 (heave) & $Z$ & $w$ & $z$\\
    4 & 횡경사 (roll)  & $K$ & $p$ & $\phi$\\
    5 & 종경사 (pitch) & $M$ & $q$ & $\theta$\\
    6 & 선수각 (yaw)   & $N$ & $r$ & $\psi$\\
  \end{tabular}};

\node[anchor=north west, align=left, text width=78mm, font=\scriptsize] at (0,-6) {%
  \textbf{직선 화살표} = 힘. \textbf{곡선 화살표} = 모멘트이며, 자기 축에 수직인
  평면 위의 원으로 그린다 (오른손 법칙). 점선은 각 원의 중심을 자기 축에 잇는 것으로,
  중심이 축에서 벗어난 타원은 이 그림이 틀리는 전형적인 방식이다.\\[4pt]
  수상선은 보통 \textbf{1 · 2 · 6}(전후 · 좌우 · 선수각)으로 줄여 다룬다.
  나머지 셋은 제어 대상에서 뺄 뿐 사라지지 않는다 — 파도가 커지면
  Roll · Pitch 가 커지고 그만큼 측정값에 잡음이 섞인다.};
\end{tikzpicture}

\end{document}